% from ega1-10-fr.tex
\oldpage[I]{184}
\begin{env}[10.3.2]
\label{I.10.3.2-fr}
Avec les notations de (\hyperref[I.10.3.1-fr]{10.3.1}), pour tout $f\in A$,
l'application canonique de
$A_{\{f\}}=\Gamma(\mathfrak{D}(f),\mathscr{O}_{\mathfrak{X}})$ dans
$\Gamma\bigl(\mathfrak{D}(f),
(\widetilde{A}/\widetilde{\mathfrak{J}})|\mathfrak{X}\bigr)
=S_f^{-1}A/S_f^{-1}\mathfrak{J}$
est \emph{surjective} et a pour noyau
$\Gamma(\mathfrak{D}(f),\mathfrak{J}^\Delta)=\mathfrak{J}_{\{f\}}$
(\hyperref[0.7.6.9-fr]{0, 7.6.9})~; ces applications définissent donc un
homomorphisme \emph{surjectif} continu, dit \emph{canonique}, du faisceau
d'anneaux topologiques $\mathscr{O}_{\mathfrak{X}}$ sur le faisceau
d'anneaux discrets
$(\widetilde{A}/\widetilde{\mathfrak{J}})|\mathfrak{X}$, dont le noyau est
$\mathfrak{J}^\Delta$~; cet homomorphisme n'est autre d'ailleurs que
$\widetilde{\varphi}$ (\hyperref[I.10.2.1-fr]{10.2.1}), où $\varphi$ est
l'homomorphisme continu $A\to A/\mathfrak{J}$~; le morphisme
$({}^a\varphi,\widetilde{\varphi}):
\operatorname{Spec}(A/\mathfrak{J})\to\mathfrak{X}$ de schémas formels
affines (où ${}^a\varphi$ est d'ailleurs l'homéomorphisme identique de
$\mathfrak{X}$ sur lui-même) est encore dit \emph{canonique}. On a donc,
d'après ce qui précède, un \emph{isomorphisme canonique}
\[
  \mathscr{O}_{\mathfrak{X}}/\mathfrak{J}^\Delta
  \xrightarrow{\sim}
  (\widetilde{A}/\widetilde{\mathfrak{J}})|\mathfrak{X}.
  \tag{10.3.2.1}\label{I.10.3.2.1-fr}
\]

Il est clair (en vertu de
$\Gamma(\mathfrak{X},\mathfrak{J}^\Delta)=\mathfrak{J}$) que
l'application $\mathfrak{J}\to\mathfrak{J}^\Delta$ est
\emph{strictement croissante}~; d'après ce qui précède, pour
$\mathfrak{J}\subset\mathfrak{J}'$, le faisceau
${\mathfrak{J}'}^\Delta/\mathfrak{J}^\Delta$ est canoniquement isomorphe à
$\widetilde{\mathfrak{J}'}/\widetilde{\mathfrak{J}}
=\widetilde{\mathfrak{J}'/\mathfrak{J}}$.
\end{env}

\begin{env}[10.3.3]
\label{I.10.3.3-fr}
Les hypothèses et notations étant toujours celles de
(\hyperref[I.10.3.1-fr]{10.3.1}), nous dirons qu'un faisceau d'idéaux
$\mathscr{J}$ de $\mathscr{O}_{\mathfrak{X}}$ est un
\emph{faisceau d'idéaux de définition de $\mathfrak{X}$} (ou un
\emph{Idéal de définition de $\mathfrak{X}$}) si, pour tout
$x\in\mathfrak{X}$, il existe un voisinage ouvert de $x$ de la forme
$\mathfrak{D}(f)$, où $f\in A$, tel que
$\mathscr{J}|\mathfrak{D}(f)$ soit de la forme $\mathfrak{H}^\Delta$, où
$\mathfrak{H}$ est un idéal de définition de $A_{\{f\}}$.
\end{env}

\begin{proposition}[10.3.4]
\label{I.10.3.4-fr}
Pour tout $f\in A$, tout Idéal de définition de $\mathfrak{X}$ induit un
Idéal de définition de $\mathfrak{D}(f)$.
\end{proposition}

\begin{proof}
Cela résulte de (\hyperref[I.10.3.1.1-fr]{10.3.1.1}).
\end{proof}

\begin{proposition}[10.3.5]
\label{I.10.3.5-fr}
Si $A$ est un anneau admissible, tout Idéal de définition de
$\mathfrak{X}=\operatorname{Spf}(A)$ est de la forme
$\mathfrak{J}^\Delta$, où $\mathfrak{J}$ est un idéal de définition de
$A$, uniquement déterminé.
\end{proposition}

\begin{proof}
En effet, soit $\mathscr{J}$ un Idéal de définition de $\mathfrak{X}$~;
par hypothèse, et puisque $\mathfrak{X}$ est quasi-compact, il y a un nombre
fini d'éléments $f_i\in A$ tels que les $\mathfrak{D}(f_i)$ recouvrent
$\mathfrak{X}$ et que
$\mathscr{J}|\mathfrak{D}(f_i)=\mathfrak{H}_i^\Delta$, où
$\mathfrak{H}_i$ est un idéal de définition de $A_{\{f_i\}}$. Pour tout
$i$, il existe donc un idéal ouvert $\mathfrak{K}_i$ de $A$ tel que
$(\mathfrak{K}_i)_{\{f_i\}}=\mathfrak{H}_i$
(\hyperref[0.7.6.9-fr]{0, 7.6.9})~; soit $\mathfrak{K}$ un idéal de
définition de $A$ contenu dans tous les $\mathfrak{K}_i$. L'image canonique
de $\mathscr{J}/\mathfrak{K}^\Delta$ dans le faisceau structural
$\widetilde{A/\mathfrak{K}}$ de
$\operatorname{Spec}(A/\mathfrak{K})$
(\hyperref[I.10.3.2-fr]{10.3.2}) est donc telle que sa restriction à
$\mathfrak{D}(f_i)$ soit égale à celle de
$\widetilde{\mathfrak{K}_i/\mathfrak{K}}$~; on en conclut que cette image
canonique est un faisceau quasi-cohérent sur
$\operatorname{Spec}(A/\mathfrak{K})$, donc de la forme
$\widetilde{\mathfrak{J}/\mathfrak{K}}$, où $\mathfrak{J}$ est un idéal de
$A$ contenant $\mathfrak{K}$ (\hyperref[I.1.4.1-fr]{1.4.1}) d'où
$\mathscr{J}=\mathfrak{J}^\Delta$
(\hyperref[I.10.3.2-fr]{10.3.2})~; en outre, comme pour tout $i$ il existe
un entier $n_i$ tel que
$\mathfrak{H}_i^{n_i}\subset\mathfrak{K}_{\{f_i\}}$, on aura, en désignant
par $n$ le plus grand des $n_i$,
$(\mathscr{J}/\mathfrak{K}^\Delta)^n=0$, et par suite
(\hyperref[I.10.3.2-fr]{10.3.2})
$(\widetilde{\mathfrak{J}/\mathfrak{K}})^n=0$, d'où finalement
$(\mathfrak{J}/\mathfrak{K})^n=0$
(\hyperref[I.1.3.13-fr]{1.3.13}), ce qui prouve que $\mathfrak{J}$ est un
idéal de définition de $A$ (\hyperref[0.7.1.4-fr]{0, 7.1.4}).
\end{proof}

\begin{proposition}[10.3.6]
\label{I.10.3.6-fr}
Soient $A$ un anneau adique, $\mathfrak{J}$ un idéal de définition de $A$
tel que $\mathfrak{J}/\mathfrak{J}^2$ soit un
$A/\mathfrak{J}$-module de type fini. Pour tout entier $n>0$, on a alors
$(\mathfrak{J}^\Delta)^n=(\mathfrak{J}^n)^\Delta$.
\end{proposition}

\begin{proof}
En effet, pour tout $f\in A$, on a (puisque $\mathfrak{J}^n$ est un idéal
ouvert)
\[
  \bigl(\Gamma(\mathfrak{D}(f),\mathfrak{J}^\Delta)\bigr)^n
  =(\mathfrak{J}_{\{f\}})^n
  =(\mathfrak{J}^n)_{\{f\}}
  =\Gamma\bigl(\mathfrak{D}(f^n),(\mathfrak{J}^n)^\Delta\bigr)
\]
